
\documentclass[../../../physics-standard_questions_all.tex]{subfiles}

\begin{document}

\setlength{\abovedisplayskip}{0mm}
\setlength{\belowdisplayskip}{0mm}

真空中に電場が在り，$xy$平面上での電位分布が
\[ \phi=a(x^2-y^2)\,(a=10\,\text{V/}\text{m}^2) \]
で与えられるものとする．

\begin{enumerate}
    
    \item $xy$平面上の$x>0$かつ$y>0$の領域について，
    $\phi=0$\,V，$\phi=\pm10$\,V，$\phi=\pm20$\,Vの
    各々に対応する等電位線の概形を図示せよ．
    また，同じ図上に点$(1\,\text{m}, \hspace{1.2mm}1\,\text{m})$
    を通る電気力線の概形を描け．

    \item 点$(1\,\text{m}, \hspace{1.2mm}0\,\text{m})$での
    電場の強さと向きを求めよ．

    \item $y$軸に沿ってなめらかに動ける正の点電荷P（質量$m$，電荷$q$）に対し，
    点$(0, \hspace{1.2mm}-h)\,(h>0)$で初速度（$y$軸の正の向きに大きさ$v_0$）を与える．
    Pが$y$軸正方向の無限遠へ飛び去るための$v_0$の条件を求めよ
    （この設問のみ文字式で答えよ）．

\end{enumerate}

\end{document}


